\section{Further Language Features}
\label{sec:scala-misc}

%% We describe pattern matching (Scala box~\ref{sb:pattern-matching}), anonymous
%% functions (Scala box~\ref{sb:anon-function}), |for| expressions (Scala
%% box~\ref{sb:for-expressions}), and how to catch
%% exceptions~\ref{sb:try-catch}). 

%% \begin{scalaBox}{Pattern matching}
%% \label{sb:pattern-matching}
The Scala operator |match| performs \emph{pattern matching}\index{pattern
matching}.  The command\ScalaIndex{match}\ScalaIndex{case}
%
\begin{scala}
  £$v$£ match{
    case £$pat_1$£ => £$cmd_1$£
    ...
    case £$pat_n$£ => £$cmd_n$£
  }
\end{scala}
tries to match the value~$v$ against each of the patterns $pat_1, \ldots,
pat_n$.  When it finds a pattern that matches, it binds any names in the
pattern against the corresponding subvalue in~$v$, and executes the corresponding
command~$cmd_i$; if several match, it executes the first such command.  If no
pattern matches, then it throws an exception.

Patterns include the following.
\begin{itemize}
\item The constant pattern, matches just that value; for example |"get"| or
  |36|.

\item A variable pattern such as ``|x|'' matches all values, and binds the
name~|x| to the value.  A typed pattern such as ``|x: A|'' matches all values
of type~|A|, and binds the name~|x| to the value.  The wildcard pattern
\SCALA{\_} matches all values.  

\item A pair pattern, e.g.~|(pat|$_1$|, pat|$_2$|)|, where |pat|$_1$ and
  |pat|$_2$ are patterns, matches a pair~|(v|$_1$\SCALA{, v}$_2$|)| provided
  $\sm v_1$ and $\sm v_2$ match |pat|$_1$ and |pat|$_2$, respectively.  Names
  are bound according to those patterns.  This extends to larger tuples in the
  obvious way.
\end{itemize}
%
Case classes (Section~\ref{sec:scala-encapsulation}), such as the |Option| type
(Section~\ref{sec:scala-types}), also allow for pattern matching.
%% \end{scalaBox}

%%%%%

%% \begin{scalaBox}{Anonymous functions}
%% \label{sb:anon-function}
An \emph{anonymous function}\index{anonymous} can be defined in Scala in a
couple of ways.  The notation \SCALA{(x: A) => exp} represents a function that
takes an argument~|x| of type~|A|, and returns the value of~|exp|.  For
example, \SCALA{(x: Int) => x+1} is a function that increments an |Int|.  The
typing ``\SCALA{: A}'' can often be omitted, when the Scala compiler can
deduce the type; for example, in \SCALA{List(1,2,3).map(x => x+1)}, the
compiler can deduce that |x| must be an |Int|.

An anonymous functions can also be defined using an underscore, ``\SCALA{_}'' to
represent a ``hole'' into which the argument is put.  For example,
\SCALA{((_:Int)+1)}  increments an |Int|.  Again, the type of the
argument can often be omitted. 
%% \end{scalaBox}

%%%%%

%% \begin{scalaBox}{\protect\SCALAKW{for} expressions}
%% \label{sb:for-expressions}
A \emph{\SCALA{for} expression}\ScalaIndexA{for expression@\protect\SCALA{for} expression}
provides a convenient way of building sequences.  The simplest form of |for|
expression is of the form
\begin{scala}
  for(x <- xs) yield f(x)
\end{scala}
where |xs| is a sequence, and |f| is a suitable function: this will generate a
sequence that contains the value |f(x)| for each |x| in~|xs|.  The ``\SCALA{x
  <- xs}'' is known as a \emph{generator}.  For example, 
\begin{scala}
  for(x <- List(1,2,3)) yield x*x
\end{scala}
will produce the list |List(1,4,9)|.

A |for| expression may contain more than one generator.  For example
\begin{scala}
  for(x <- List(1,2,3); y <- List(4,5)) yield x*y
\end{scala}
will produce |List(4,5,8,10,12,15)|.  
%% Note that the |for| expression iterates
%% over the second generator for each value from the first generator, before
%% moving on to the next value from the first generator.

A |for| expression may also contain a \emph{filter}, restricting the range of
values from which results are produced.  For example
\begin{scala}
  for(x <- List(1,2,3,4); if x%2 == 0; y <- List(1,2,3)) yield x*y
\end{scala}
only calculates |x*y| for even values of~|x|, so produces
|List(2, 4, 6, 4, 8, 12)|.

%% Scala has several different types of sequence: we mostly use |List|s in this
%% book.  The type of the result produced by a |for| expression depends upon the
%% type of the sequences used in the generators.  An arbitrary sequence can be
%% converted into a |List| using the operation |toList|, for example
%% \begin{scala}
%%   (for(x <- 1 to 3) yield x*x).toList
%% \end{scala}
%% \end{scalaBox}


%% \begin{scalaBox}{Catching exceptions}
%% \label{sb:try-catch}
Exceptions can be caught using a |try ... catch|%
\ScalaIndex{try}\ScalaIndex{catch} construct as follows.
\begin{scala}
try{ prog } catch{ handler }
\end{scala}
%
Here |prog| is a program fragment, and |handler| is a partial function
describing how to handle any exception thrown by |prog|.  The exception
handler is normally written using pattern matching syntax (described above),
for example%
\begin{scala}
  try{ prog } catch{
    case st: Stopped => ...
    case ae: AssertionError => ...
  }
\end{scala}
%% \end{scalaBox}


%% \begin{scalaBox}{{\scalashape apply} functions}
%% \label{sb:apply}
If an object~|o| has an operation called |apply|\ScalaIndex{apply} as follows:
\begin{scala}
  def apply(x: A) = ...
\end{scala}
then it can be invoked as |o(e)| where  |e| is an expression of type~|A|.
This generalises to different numbers of parameters, including none.  
%% \end{scalaBox}

\framebox{TODO:} |import|,  equality and hash codes
(from datatypes3).
